\documentclass{article}
\usepackage[margin=1in]{geometry}
\usepackage{tabularx}
\usepackage{array}
\usepackage{booktabs}
\usepackage{xcolor}
\usepackage{amsmath}
\usepackage[T1]{fontenc}
\usepackage[spanish,es-nodecimaldot]{babel}
\usepackage{graphicx,enumitem,fancyhdr,listings,hyperref}
\definecolor{azul}{HTML}{1B2A5E}
\definecolor{dorado}{HTML}{E5B300}
\definecolor{fondo}{HTML}{F3F6FA}
\hypersetup{colorlinks=true,urlcolor=azul,pdftitle={Ejercicios Econometría Taller 1 - Introducción a la computación estadística y econometría},pdfauthor={Iván Canales}}
\pagestyle{fancy}\fancyhf{}
\fancyhead[L]{\small\color{azul}Econometría - Taller 1}
\setlength{\headheight}{32pt}
\fancyhead[R]{\includegraphics[height=0.95cm]{../Logo_UTFSM.png}}
\fancyfoot[C]{\small\thepage}
\renewcommand{\headrulewidth}{0.4pt}
\setlength{\parindent}{0pt}\setlength{\parskip}{5pt}
\setlist[enumerate]{leftmargin=*,itemsep=13pt,parsep=4pt}
\lstset{language=R,basicstyle=\ttfamily\small,backgroundcolor=\color{fondo},frame=single,rulecolor=\color{fondo},breaklines=true,columns=fullflexible,keepspaces=true,showstringspaces=false,xleftmargin=6pt,xrightmargin=6pt,aboveskip=6pt,belowskip=6pt}
\newcommand{\codigo}[1]{\texttt{\detokenize{#1}}}
\newcommand{\captura}[3]{\begin{minipage}[t]{0.48\linewidth}\centering\includegraphics[width=\linewidth,height=#3,keepaspectratio]{colab-#1.png}\par\vspace{2pt}{\footnotesize\color{azul}\textbf{#1.} #2}\end{minipage}}
\begin{document}
{\LARGE\bfseries\color{azul}Ejercicios Econometría Taller 1\par}
{\large\color{azul}Introducción a la computación estadística y econometría\par}
\vspace{2pt}{\color{dorado}\hrule height 1.5pt}\vspace{7pt}
{\large\bfseries\color{azul}Si no tienes aún Google Colaboratory, esta es la manera de hacerlo. La otra opción es buscarlo directamente en Google.\par}
{\small Accede con tu cuenta de Google. Si Colaboratory no aparece en Drive, usa \textbf{Nuevo $\rightarrow$ Más $\rightarrow$ Conectar más aplicaciones}, busca la aplicación oficial y conéctala. Luego abre el notebook y selecciona el entorno \textbf{R}.}

\captura{1}{Abre Drive desde las aplicaciones de Google.}{4.3cm}\hfill
\captura{2}{Accede a Google Workspace Marketplace.}{4.3cm}
\par\vspace{8pt}
\captura{3}{Busca Colaboratory y pulsa Instalar.}{3.9cm}\hfill
\captura{4}{En Drive, entra a Nuevo $\rightarrow$ Más $\rightarrow$ Google Colaboratory.}{3.9cm}
\par\vspace{8pt}
\captura{5}{Entorno de ejecución $\rightarrow$ Cambiar tipo de entorno de ejecución.}{4.5cm}\hfill
\captura{6}{Selecciona R y pulsa Guardar.}{4.5cm}
\par\vspace{5pt}
{\footnotesize Con R seleccionado, ejecuta el código directamente: no necesitas \texttt{\%\%R} ni cargar la extensión de Python \codigo{rpy2}.}

\newpage
{\Large\bfseries\color{azul}Repaso del taller 1\par}
\vspace{2pt}{\color{dorado}\hrule height 1pt}\vspace{6pt}
\textbf{A continuación se hará un repaso de lo visto en el taller 1.}

Haz un código en Colab explicando lo que hiciste. Para ello, averigua cómo crear una celda Markdown.

\begin{enumerate}
\item \textbf{Nombre y apellido con ciclos.} Escribe un programa que muestre tu nombre y tu apellido de manera alternada: nombre, apellido, nombre, apellido, etc. Elige e indica cuántas alternancias realizará. Resuélvelo primero con \codigo{while} y después con \codigo{for}.

\textit{Consejo: utiliza las sentencias condicionales \codigo{if} y \codigo{else} para decidir qué texto mostrar en cada iteración.}

\item \textbf{Una sucesión de signos alternados.} La sucesión
\[
1,\,-1,\,2,\,-2,\,3,\,-3,\,4,\,-4,\,\ldots
\]
tiene una estructura que puede codificarse. Genérala hasta llegar a $-20$, inclusive, utilizando primero \codigo{while} y después \codigo{for}.

\textit{Consejo: utiliza \codigo{if} y \codigo{else} para alternar el signo.}

\item \textbf{Funciones para construir secuencias y vectores.} Explica para qué sirven \codigo{c()}, \codigo{seq()}, \codigo{numeric()} y \codigo{rep()}. Investiga también el uso de \codigo{seq_len()}, \codigo{seq_along()}, \codigo{sequence()} y \codigo{gl()}. Para cada función, presenta un ejemplo en R y explica su resultado.

En particular, ejecuta \codigo{gl(2, 3, labels = c("A", "B"))} y explica para qué sirve cada dato que escribes entre los paréntesis y cómo se organiza el resultado.

\item \textbf{Valores aleatorios y redondeo.} La función \codigo{runif()} permite generar números aleatorios de una distribución uniforme; por defecto, sus límites son 0 y 1. Investiga sus parámetros y crea un vector de longitud 1.000. Asígnalo mediante \codigo{<-} a un objeto llamado \codigo{v}.

Luego ejecuta \codigo{floor(v * 100)}. ¿Qué hace \codigo{floor()} y qué valores puede tomar el resultado? Compara algunos elementos de \codigo{v} con sus correspondientes valores transformados.
\end{enumerate}

\newpage
{\Large\bfseries\color{azul}Aplicación con datos de Wooldridge\par}
\vspace{2pt}{\color{dorado}\hrule height 1pt}\vspace{6pt}
\begin{enumerate}[start=5]
\item \textbf{Carga de datos y transformación logarítmica.} Recupera la base utilizada en clase. Si todavía no tienes instalado el paquete, ejecuta la instalación una vez; después carga los datos:
\begin{lstlisting}
install.packages("wooldridge", repos = "https://cloud.r-project.org/")
datos <- wooldridge::wage1
\end{lstlisting}
Asigna \codigo{y <- datos$wage} y ejecuta \codigo{hist(log(y))}. Agrega un título al gráfico y nombres a sus ejes. Explica con tus palabras qué observas. Recuerda que \codigo{wage} es el salario por hora, expresado en USD de 1976, y que \codigo{log()} calcula el logaritmo natural.

\item \textbf{Antigüedad y gráficos de caja.} Crea una variable nueva en \codigo{datos} que clasifique \codigo{tenure} en cuatro grupos: 0 años; entre 1 y 3 años, inclusive; entre 4 y 8 años, inclusive; y más de 8 años. \codigo{tenure} mide la antigüedad con el empleador actual, en años.

Construye gráficos de caja del salario \codigo{wage} según esta nueva variable. Compara los grupos y explica lo que observas en sus medianas, dispersión y posibles valores atípicos.

\item \textbf{Distribución de promedios muestrales.} Crea un vector llamado \codigo{avg_wage} de longitud 1.000. Toma una muestra aleatoria simple \textbf{sin reposición} de 50 filas de \codigo{datos}, calcula el promedio de \codigo{wage} con \codigo{mean()} y almacénalo en la posición 1 de \codigo{avg_wage}.

Vuelve a tomar una muestra de 50 filas de la base completa, calcula su promedio salarial y guárdalo en la posición 2. Continúa mediante un ciclo hasta completar las 1.000 posiciones.

La selección es sin reposición \textbf{dentro de cada muestra}; las muestras se toman nuevamente desde la base completa, por lo que una persona puede aparecer en distintas repeticiones.

Antes de iniciar el ciclo, fija la siguiente semilla una sola vez:
\begin{lstlisting}
set.seed(123)
\end{lstlisting}
Dentro del ciclo, utiliza esta sentencia para seleccionar las 50 filas de cada muestra:
\begin{lstlisting}
filas <- sample(1:nrow(datos), size = 50, replace = FALSE)
\end{lstlisting}
Usa \codigo{datos[filas, ]} para acceder a las filas seleccionadas.

Finalmente, construye un histograma de \codigo{avg_wage} con \codigo{hist()}. ¿Qué forma observas? Relaciona lo observado con la distribución de los promedios muestrales.

\item \textbf{Nube de puntos y correlación.} Construye un diagrama de dispersión entre \codigo{educ} y \codigo{exper}, y calcula su correlación con \codigo{cor()}. Ambas variables se miden en años: educación y experiencia laboral, respectivamente.

Explica qué observas en la nube de puntos y cómo se relaciona con el signo y la magnitud de la correlación. Distingue asociación de causalidad.
\end{enumerate}
\end{document}
